\section{Global finite-order radial TE control with bounded positive index}
\label{cc:section}

The following construction turns the angular-momentum term in the
actual radial Maxwell equation into a controllable finite ensemble.
It uses only finite concentric layers with
$1\leq\varepsilon\leq q^2<\varepsilon_+$. Its support may grow.
The result is approximate global transfer control for every
prescribed finite set of TE orders. It is not a bypass theorem.

\subsection{An identity loop with an amplified angular derivative}

Use $x=k_0r$, the electric TE state $v=(u,u')^T$, and
\begin{equation}\label{cc:matrices}
 X=\begin{pmatrix}0&1\\0&0\end{pmatrix},\qquad
 Y=\begin{pmatrix}0&0\\1&0\end{pmatrix},\qquad J=X-Y.
\end{equation}
For $\nu=l(l+1)$ the exact equation in an index-$n$ layer is
\begin{equation}\label{cc:actual-ode}
 v'=(X-n^2Y+\nu x^{-2}Y)v .
\end{equation}
The state is continuous at dielectric interfaces. No magnetic or
surface parameter is being varied.

Fix $1<q<\sqrt{\varepsilon_+}$. A high quarter layer has index $q$
and length $\pi/(2q)$; a low quarter layer has index one and length
$\pi/2$. The reference equation omits $\nu x^{-2}Y$. Its quarter
matrix is
$T_n=\left(\begin{smallmatrix}0&1/n\\-n&0\end{smallmatrix}\right)$.
A loop consists of $M$ high--low pairs followed by $M$ low--high
pairs. Its length and reference endpoint are
\begin{equation}\label{cc:loop-length}
 \ell_M=M\pi(1+q^{-1}),\qquad U_0(\ell_M)=I.
\end{equation}
Indeed, the first half gives
$(-1)^M\operatorname{diag}(q^M,q^{-M})$ and the second its inverse.

\begin{lemma}[PROVED: exact angular derivative of a finite dielectric loop]
\label{cc:loop-derivative}
Let $B_M=\int_0^{\ell_M}U_0(s)^{-1}YU_0(s)\,ds$. Then
\begin{align}
 B_M&=\begin{pmatrix}0&-d_M\\b_M&0\end{pmatrix},\notag\\
 b_M&=\frac{\pi(q^2+q^{-1})}{2(q^2-1)}(q^{2M}-1),\notag\\
 d_M&=\frac{\pi(1+q^{-1})}{2(q^2-1)}(1-q^{-2M}),\notag\\
 \frac{d_M}{b_M}&=\kappa_q q^{-2M},\qquad
 \kappa_q=\frac1{q^2-q+1}<1. \label{cc:loop-derivative-formula}
\end{align}
Moreover
\begin{equation}\label{cc:integrated-norm}
 \int_0^{\ell_M}\|U_0(s)^{-1}YU_0(s)\|\,ds
       =b_M+d_M\leq2b_M .
\end{equation}
\end{lemma}
\begin{proof}
If the first row of $U_0$ is $(a,b)$, determinant one gives
\begin{equation}\label{cc:conjugated-y}
 U_0^{-1}YU_0=
 \begin{pmatrix}-ab&-b^2\\a^2&ab\end{pmatrix}.
\end{equation}
This rank-one matrix has norm $a^2+b^2$. In high and low layers
use $0\leq\theta\leq\pi/2$, with $ds=d\theta/q$ and
$ds=d\theta$, respectively. Up to an irrelevant common sign,
the first-row pairs in the forward half, for $0\leq j<M$, are
\begin{equation}\label{cc:forward-rows}
 (q^j\cos\theta,q^{-j-1}\sin\theta),\qquad
 (-q^{j+1}\sin\theta,q^{-j-1}\cos\theta).
\end{equation}
Those in the reverse half, for $M\geq h\geq1$, are
\begin{equation}\label{cc:reverse-rows}
 (q^h\cos\theta,q^{-h}\sin\theta),\qquad
 (-q^{h-1}\sin\theta,q^{-h}\cos\theta).
\end{equation}
Integrating $a^2$ over these layers gives
$b_M=\frac{\pi}{2}(q^{-1}+q^2)\sum_{j=0}^{M-1}q^{2j}$.
Integrating $b^2$ gives
$d_M=\frac{\pi}{2}(q^{-3}+q^{-2})\sum_{j=0}^{M-1}q^{-2j}$.
The integral of $ab$ in each forward pair is
$(q^{-2}-1)/2$, and in each reverse pair its negative.
The diagonal entries therefore vanish exactly. Summing the
geometric series gives \eqref{cc:loop-derivative-formula}.
The rank-one norm identity gives \eqref{cc:integrated-norm}.
\end{proof}

\begin{lemma}[PROVED: finite-radius error for one actual Maxwell loop]
\label{cc:one-loop}
Let $U_\nu(r)$ be the actual transfer through this loop starting at
dimensionless radius $r>0$, and let $0\leq\nu\leq V$. Then
\begin{equation}\label{cc:one-loop-error}
 \left\|U_\nu(r)-I-\frac{\nu}{r^2}B_M\right\|
 \leq \frac{4\nu\ell_Mb_M}{r^3}
       +\frac{2\nu^2b_M^2}{r^4}\exp(2\nu b_M/r^2).
\end{equation}
Also $\|U_\nu(r)\|\leq\exp(2\nu b_M/r^2)$.
\end{lemma}
\begin{proof}
Writing the actual fundamental matrix as $U_0W$ gives
\[
 W'(s)=\frac{\nu}{(r+s)^2}U_0(s)^{-1}YU_0(s)W(s),
 \qquad W(0)=I.
\]
The integrated coefficient norm is at most $\beta=2\nu b_M/r^2$.
The Volterra series beyond its linear term is bounded by
$e^\beta-1-\beta\leq\beta^2e^\beta/2$.
Replacing $(r+s)^{-2}$ in that linear term by $r^{-2}$ costs
at most $2\nu\ell_M(b_M+d_M)/r^3\leq4\nu\ell_Mb_M/r^3$.
Since $U_0(\ell_M)=I$, these are the endpoint errors.
The same series bounds $\|W(\ell_M)\|$ by $e^\beta$.
\end{proof}

\subsection{A macroscopic angular shear in a thin distant annulus}

\begin{theorem}[PROVED: simultaneous angular shear from positive-index layers]
\label{cc:angular-shear}
Fix finitely many orders $l_1,\ldots,l_m\geq1$, put
$\nu_j=l_j(l_j+1)$ and $V=\max_j\nu_j$, and fix $t\geq0$.
For every sufficiently large starting radius $R$ there are finite
layers of index one or $q$ in an annulus of dimensionless thickness
$o(R)$ whose actual TE transfers converge simultaneously to
\begin{equation}\label{cc:angular-shear-target}
 \exp(t\nu_jY)=I+t\nu_jY,\qquad 1\leq j\leq m .
\end{equation}
The material bound is independent of $R$, $t$, and the orders.
\end{theorem}
\begin{proof}
For $t=0$ use no layers. Otherwise fix $a\in(1/2,1)$ and set
\begin{equation}\label{cc:scales}
 M=\lfloor a\log_qR\rfloor,\quad h=b_M/R^2,\quad
 K=\lfloor t/h\rfloor
\end{equation}
for sufficiently large $R$. Use $K$ consecutive loops, starting
at $r_i=R+i\ell_M$, $0\leq i<K$. The explicit formula for $b_M$
gives
\[
 h\to0,\qquad K\ell_M/R=O(R^{1-2a}\log R)\to0,\qquad
 q^{-2M}\to0.
\]
Set $\epsilon_M=\kappa_qq^{-2M}$. One actual loop differs from
$I+h\nu Y$ by at most
\begin{equation}\label{cc:step-error}
 \nu h\epsilon_M+\frac{2\nu hK\ell_M}{R}
       +\frac{4\nu h\ell_M}{R}+2\nu^2h^2e^{2\nu h}.
\end{equation}
The first two terms compare $(\nu/r_i^2)B_M$ with $h\nu Y$,
using $B_M/b_M=Y-\epsilon_MX$ and
$|R^2/r_i^2-1|\leq2i\ell_M/R$. The last two terms are
\eqref{cc:one-loop-error}.

Each actual step has norm at most $e^{2Vh}$ and each ideal step
at most $e^{Vh}$. Telescoping the $K$-fold products bounds their
difference by $Ke^{2Vt}$ times \eqref{cc:step-error} with $\nu=V$.
Since $Y^2=0$, the ideal product is $I+Kh\nu Y$, whose difference
from the target is at most $Vh$. Thus the simultaneous endpoint
error is at most
\begin{equation}\label{cc:global-error}
 e^{2Vt}t\left[
 V\epsilon_M+\frac{2VK\ell_M}{R}
       +\frac{4V\ell_M}{R}+2V^2h e^{2Vh}\right]+Vh.
\end{equation}
Every term tends to zero. The total thickness is $K\ell_M=o(R)$.
Every layer has real permittivity one or $q^2$, and every
structure at finite $R$ has finitely many interfaces.
\end{proof}

\subsection{Finite-ensemble global transfer control}

\begin{definition}\label{cc:realizable}
For the fixed finite orders, a tuple
$(G_1,\ldots,G_m)\in{\rm SL}(2,\mathbb R)^m$ is
\emph{asymptotically realizable} if, for every sufficiently large
$R$, a finite radial dielectric in $[R,R+d_R]$, with
$d_R/R\to0$ and $1\leq\varepsilon\leq q^2$, has simultaneous
TE transfers converging to that tuple as $R\to\infty$.
\end{definition}

\begin{lemma}[PROVED: common real symplectic control and closure]
\label{cc:common-and-closure}
Every common tuple $(G,\ldots,G)$, $G\in{\rm SL}(2,\mathbb R)$,
is asymptotically realizable. The realizable tuples form a closed
multiplicative semigroup.
\end{lemma}
\begin{proof}
In the reference equation a vacuum layer realizes $e^{\theta J}$
for any real $\theta$, by adding a multiple of $2\pi$ to make the
length nonnegative. A high--low quarter pair with adjustable high
index $1\leq s\leq q$ realizes
$-\operatorname{diag}(s,s^{-1})$; a vacuum half-period removes
the sign. Products give any positive diagonal
$\operatorname{diag}(s,s^{-1})$, $s>0$, with inverses obtained
by interchanging layer order. The singular-value decomposition
expresses every real determinant-one matrix as two proper
rotations and such a diagonal. Thus finite bounded-index layers
realize any common $G$ exactly in the reference equation.
For this fixed sequence, the interaction estimate above bounds
the angular correction by $O(R^{-2})$, uniformly over the finite
orders. Its length is fixed, proving common realizability.

Two realizing families concatenate: the first endpoint divided
by $R$ tends to one, so use the second family at that endpoint.
Total thickness remains $o(R)$ and transfers converge to the
product. For closedness let realizable tuples $G^{(n)}\to G$.
Choose increasing thresholds $R_n\to\infty$ so that the family
for $G^{(n)}$, for every $R\geq R_n$, has transfer error at most
$1/n$ and thickness at most $R/n$. On $R_n\leq R<R_{n+1}$
use that family. Transfers converge to $G$ and relative thickness
to zero. All tuple norms are maxima of the finitely many matrix
norms. The same reasoning permits limits of finite products.
\end{proof}

\begin{theorem}[PROVED: global approximate finite-order TE transfer control]
\label{cc:global-control}
For distinct positive integers $l_1,\ldots,l_m$ and arbitrary
$G_1,\ldots,G_m\in{\rm SL}(2,\mathbb R)$, their tuple is
asymptotically realizable by actual finite ordinary dielectrics,
for any fixed $1<q<\sqrt{\varepsilon_+}$. Consequently each
prescribed accuracy is attained by a finite bounded positive
scalar radial structure. No uniform radius, layer-count, or
thickness bound is asserted as target or accuracy varies.
\end{theorem}
\begin{proof}
In tuple notation put
$A=(\nu_jY)_j$, $B=(-\nu_jX)_j$, and $C=A+B=(-\nu_jJ)_j$.
Theorem~\ref{cc:angular-shear} supplies $e^{tA}$ for $t\geq0$.
Common conjugation by $J$ supplies $e^{tB}$ for $t\geq0$,
because $JYJ^{-1}=-X$. The finite-dimensional product formula
\[
 e^{tC}=\lim_{n\to\infty}(e^{tA/n}e^{tB/n})^n
\]
and semigroup closure supply $e^{tC}$ for $t\geq0$.
Every $\nu_j$ is an integer, hence $e^{2\pi C}=I$.
Reduction modulo $2\pi$ supplies $e^{tC}$ for every real $t$.
Another product formula gives
\[
 e^{-tB}=e^{t(A-C)}
       =\lim_{n\to\infty}(e^{tA/n}e^{-tC/n})^n,\qquad t\geq0.
\]
Common conjugation then supplies $e^{-tA}$. This explicitly
generates inverses without negative physical layer thicknesses.

Put $H=[X,Y]=\operatorname{diag}(1,-1)$. The Lie algebra
generated by common $X,Y,H$ and $(\nu_jY)_j$ contains
\[
 (\nu_j^pX)_j,\quad(\nu_j^pY)_j,\quad(\nu_j^pH)_j
       \qquad(p=0,1,2,\ldots).
\]
Indeed $[H,X]=2X$, $[H,Y]=-2Y$, and bracketing
$(\nu_j^pX)_j$ with $(\nu_jY)_j$ raises the power to $p+1$.
The distinct numbers $\nu_j=l_j(l_j+1)$ give an invertible
Vandermonde matrix for powers $0,\ldots,m-1$. Linear combinations
therefore isolate $X,Y,H$ in each coordinate. The generated Lie
algebra is all of $\mathfrak{sl}(2,\mathbb R)^m$.

For completeness this yields global group control as follows.
Let $\Gamma$ be the subgroup generated by the common
one-parameter groups and $e^{tA}$ for all real $t$.
The span of their generators conjugated by $\Gamma$ is invariant
under their adjoint flows. Differentiating these flows shows it
contains the generated Lie algebra, hence the full algebra just
computed. Select $3m$ conjugated generators forming a basis.
The product of their one-parameter groups has invertible
derivative at the identity, so the inverse function theorem puts
an open identity neighborhood in $\Gamma$. It is an open
subgroup. The singular-value decomposition also shows
$\mathrm{SL}(2,\mathbb R)^m$ is connected, so $\Gamma$ is the
whole group. Every generator and inverse belongs to the closed
realizable semigroup, so all of $\Gamma$ does. This proves the
claim using actual radial Maxwell transfer matrices.
\end{proof}

\begin{corollary}[PROVED: prescribed finite TE core gains and relative phases]
\label{cc:core-gains}
Fix finitely many distinct TE orders and arbitrary nonzero complex
numbers $g_j$. There is a sequence of the finite ordinary radial
dielectrics above, with vacuum-core radius tending to infinity,
whose physical core multipliers converge to $g_j$ up to one
common unit-modulus factor. This includes simultaneous large gains
for both electric inversion parities. After removing that common
phase, convergence of the fields and their derivatives is uniform
on every fixed compact subset of the vacuum core, for every
coherent input in the chosen finite angular band.
\end{corollary}
\begin{proof}
Restrict the starting radius to $R=2\pi n\to\infty$.
The origin-normalized regular solution is $\psi_l(x)=xj_l(x)$;
its state at $R$ tends to the unit vector
\[
 v_l^0=(\sin(-l\pi/2),\cos(-l\pi/2))^T .
\]
For each order choose a nonzero real vector $w_j$ satisfying
$w_{j,1}+iw_{j,2}=i^{l_j+1}/g_j$.
There is a real determinant-one matrix taking $v_{l_j}^0$
to $w_j$, for example
\[
 G_j=
 \left[w_j,\frac{-Jw_j}{|w_j|^2}\right]
 \left[v_{l_j}^0,-Jv_{l_j}^0\right]^{-1}.
\]
Apply Theorem~\ref{cc:global-control} to these fixed targets.
The actual regular state at the outer endpoint $X_R=R+d_R$
then converges to $w_j$.

The exterior incoming coefficient is
$A_j=(u\xi_{l_j}^{+\prime}-u'\xi_{l_j}^+)/(2i)$.
The finite Hankel formula at $X_R\to\infty$ gives
\[
 \xi_l^+(X_R)=(-i)^{l+1}e^{iX_R}(1+o(1)),\qquad
 \xi_l^{+\prime}(X_R)=i(-i)^{l+1}e^{iX_R}(1+o(1)).
\]
Consequently
$e^{-iX_R}A_j\to1/(2g_j)$.
A free regular wave has incoming coefficient $1/2$, so the
physical relative core multiplier is $(2A_j)^{-1}$.
Multiplying that multiplier by the common phase $e^{iX_R}$
gives the asserted limit $g_j$. The selected band is finite and
its regular vacuum basis is smooth on each compact set, which
proves uniform coherent field and derivative convergence there.
All azimuthal indices within an order have the same multiplier.
\end{proof}

\begin{remark}[Scope and the remaining channel requirement]
\label{cc:scope}
Theorem~\ref{cc:global-control} overcomes the fixed-annulus ceiling
by proving global approximate reachability when support may grow.
It neither assumes a circuit model nor infers global reachability
from the earlier local inverse theorem. Only TE and finitely many
prescribed orders are controlled; all azimuthal indices of an order
share its transfer. Corollary~\ref{cc:core-gains} shows that the
reciprocal-parity gain pattern of the common stack is not a
restriction on all bounded positive radial structures.

Transfer targets alone do not supply counted channels. A proposed
subspace must additionally satisfy positive passage and absolute
whole-plane bypass for every coherent input, in the physical input
power norm. Fields in the many concentric layers can be large.
No estimate of their normal plane flux follows from endpoint
control. Neither a bypass estimate nor a qualifying many-channel
subspace is asserted here.
\end{remark}
